\section{Motivation}
\label{sec:motivation}

\subsection{Who cares about TTFT}

Before arguing that prefill can be moved off the critical path, it is worth
stating why the quantity being moved matters.

\paragraph{TTFT is a provisioning metric.} Serving systems are sized and evaluated
against TTFT service-level objectives. Cutting TTFT at fixed hardware means more
requests are served inside the same SLO, which is a capacity statement, not only
a user-experience one.

\paragraph{Prefill is the part of latency still addressable at the systems layer.}
Decode time is set by memory bandwidth and by how many tokens the model chooses
to emit. A serving system cannot shorten it without changing the model or the
output. Prefill, by contrast, is work whose \emph{timing} the system controls.

\paragraph{Prefill removed from the critical path is throughput, not just latency.}
Prefill and decode contend for the same GPU. Work performed during an idle window
is work the loaded engine no longer has to interleave with decode. From the
serving system's point of view the underlying quantity above TTFT is goodput, and
relocating prefill improves it in a way that trimming a latency constant does not.

\paragraph{Agentic call sites are prefill-heavy by construction.} In many agent
applications the individual LLM calls consume a large context and emit a small
structured output. For such call sites TTFT is a dominant share of the call's
latency, so it is the term worth attacking.

Concretely, three parties benefit: the application developer, through perceived
responsiveness; the operator, through SLO attainment per GPU; and single-GPU or
local deployments, where there is no spare replica to hide a cold prefill behind.

\subsection{Why a compiler can know the prompt}

Figure~\ref{fig:motivation} shows the structural opportunity. An agent program
alternates between LLM calls and non-LLM work --- tool invocations, retrieval,
environment interaction. During the non-LLM work the engine has nothing to do for
this request. That gap is not incidental; it is the defining rhythm of agentic
execution, and it recurs at every hop.

The question is what can usefully be computed inside the gap. Filling it requires
knowing the content of a prompt that has not been issued yet. For a program
written against byLLM, most of that content is already determined: the meaning-typed
abstraction constructs the prompt from typed program constructs rather than from
developer-authored strings~\cite{dantanarayana2025mtp}, so the instruction text,
the type and schema context, and the surrounding graph context are all fixed once
the program is compiled. What is released only at runtime is the set of parameter
values --- and even those are not uniformly late: many of them are computed by an
earlier statement, well before the call that consumes them is reached. A static
analysis that tracks where each argument is produced can therefore tell the
runtime not only \emph{what} the future prompt will look like but \emph{how much
of it is already available now}.

\begin{figure}[t]
\centering
% Idle-window timeline: what the engine does during a tool/RAG gap.
% Box widths are chosen so every label fits at \scriptsize; \fittext is the
% safety net that shrinks anything still too wide.
{\setlength{\fboxsep}{0pt}\setlength{\unitlength}{1mm}
\begin{picture}(84,40)(0,0)

  % ---- gap annotation
  \put(35,31){\line(0,1){3}}
  \put(56,31){\line(0,1){3}}
  \put(35,32.5){\vector(1,0){21}}
  \put(30,34.5){\makebox(31,4)[c]{\scriptsize tool call / retrieval gap}}

  % ---- lane labels
  \put(0,25){\makebox(15,6)[l]{\scriptsize\bfseries Baseline}}
  \put(0,9){\makebox(15,6)[l]{\scriptsize\bfseries \sysname}}

  % ---- baseline lane
  \lbox{16}{25}{6}{6}{cbase}{pf$_A$}
  \lbox{22}{25}{13}{6}{cdec}{decode $A$}
  \lbox{35}{25}{21}{6}{cidle}{\emph{engine idle}}
  \lbox{56}{25}{15}{6}{cbase}{prefill $B$}
  \lbox{71}{25}{13}{6}{cdec}{decode $B$}

  % TTFT bracket, baseline
  \put(56,21){\line(0,1){4}}
  \put(71,21){\line(0,1){4}}
  \put(56,22.5){\vector(1,0){15}}
  \put(48,17.5){\makebox(31,4)[c]{\scriptsize TTFT$_B$ (long)}}

  % ---- ours lane
  \lbox{16}{9}{6}{6}{cours}{pf$_A$}
  \lbox{22}{9}{13}{6}{cdec}{decode $A$}
  \lbox{35}{9}{21}{6}{cspec}{\bfseries spec.\ prefill $B$}
  \lbox{56}{9}{5}{6}{cours}{pf}
  \lbox{61}{9}{13}{6}{cdec}{decode $B$}

  % TTFT bracket, ours
  \put(56,5){\line(0,1){4}}
  \put(61,5){\line(0,1){4}}
  \put(56,6.5){\vector(1,0){5}}
  \put(43,1){\makebox(31,4)[c]{\scriptsize TTFT$_B$ (short)}}

\end{picture}}

\caption{The idle window created by a tool call or retrieval round trip. In the
baseline the engine is idle through the gap and pays a cold prefill for the next
call. \sysname{} spends the gap constructing the longest knowable prefix of that
call, so the demand-time prefill shrinks to the residual.}
\label{fig:motivation}
\end{figure}
